\documentclass[11pt,a4paper]{article}
\usepackage[T1]{fontenc}
\usepackage{lmodern,amsmath,amssymb,amsthm,mathtools,booktabs,microtype}
\usepackage[margin=25mm]{geometry}
\usepackage[colorlinks=true,linkcolor=blue,citecolor=blue,urlcolor=blue]{hyperref}
\setlength{\emergencystretch}{2em}
\newtheorem{theorem}{Theorem}[section]
\newtheorem{lemma}[theorem]{Lemma}
\newtheorem{proposition}[theorem]{Proposition}
\newcommand{\Q}{\mathbb Q}\newcommand{\Z}{\mathbb Z}
\DeclareMathOperator{\rad}{rad}\DeclareMathOperator{\Tr}{Tr}
\title{An eleven-variable construction under rational \texorpdfstring{$abc$}{abc}}
\author{}\date{17 November 2026}
\begin{document}\maketitle

\begin{abstract}
A derivative quotient recovers simple collisions between two translated
contact polynomials. Under ordinary rational $abc$, the unrecovered
multiple factors have a height cost independent of the contact degree.
Sixteen symmetric shifts give a single detector containing the input and
root-height information. A balanced transformation of the six elliptic
tests makes every new odd denominator prime harmless for completeness.
Using one two-parameter intrinsic pole block and the guarded finite-algebra
encoder gives eleven variables, including the distinguished input.
The construction uses the retained elliptic group and formal-comparison
theorems, simultaneous prime-pattern theorem, trace-height theorem, and
Burungale's simultaneous-twist theorem.
\end{abstract}

\section{Radicals of translated power values}

All heights are absolute logarithmic heights. For a nonzero rational
number $r$, set
\[
 r^+_{\mathrm{odd},S}(r)=
 \prod_{\substack{p\notin S\\v_p(r)>0\text{ odd}}}p.
\]
For a nonzero integer $N$, its simple-prime support is
$\prod_{v_p(N)=1}p$ and its multiple-prime support is
$\prod_{v_p(N)\ge2}p$. A fixed finite set $S$ changes their logarithms
by bounded constants.

We use the following standard consequence of ordinary rational $abc$:
for every fixed squarefree homogeneous form $H\in\Z[A,B]$ and every
$\epsilon>0$,
\begin{equation}\label{eq:polyabc}
 \log\rad H(A,B)\ge(\deg H-2-\epsilon)\log\max(|A|,|B|)-O(1)
\end{equation}
for primitive integer pairs with $H(A,B)\ne0$.
This is the Belyi-form consequence stated in
\cite[Appendix A.2, equation (A.2)]{RSZ}.
All forms and exponents used below are fixed before the coefficient
shift or source tuple is chosen.

\begin{lemma}\label{lem:simple}
Let $G(X)=\prod_{i=1}^k(X+c_i)$ for fixed distinct integers $c_i$.
For $G(x)\ne0$, let $T(x)$ be the product of primes outside $S$
with $v_p(G(x))=1$, and $J(x)$ the product with
$v_p(G(x))\ge2$. Then
\[
 \log T(x)\ge(k-4-2\epsilon)h(x)-O(1),\qquad
 \log J(x)\le(2+\epsilon)h(x)+O(1).
\]
\end{lemma}
\begin{proof}
Write $x=A/B$ in lowest terms, $B>0$, and
$N=\prod_i(A+c_iB)$. Then $\gcd(N,B)=1$.
The squarefree homogeneous form $\prod_i(A+c_iB)$ has degree $k$,
so \eqref{eq:polyabc} gives
$\log\rad N\ge(k-2-\epsilon)h(x)-O(1)$.
Also $\log|N|\le k h(x)+O(1)$. Prime by prime,
\[
 \log J(x)\le\log|N|-\log\rad N.
\]
Finally $\log T(x)=\log\rad N-\log J(x)+O(1)$.
\end{proof}

\begin{lemma}\label{lem:power-multiple}
Fix $m\ge1$ and distinct nonzero rationals $s_1,\ldots,s_r$.
For $M\ne0$ with $D(M):=\prod_j(2M^m+s_j)\ne0$, the product
of primes outside a fixed set with $v_p(D(M))\ge2$ has logarithm
at most
\[
                         (2+\epsilon)h(M)+O(1).
\]
The constants may depend on $m$ and the shifts, but not on $M$.
\end{lemma}
\begin{proof}
Write $M=A/B$ in lowest terms and $s_j=p_j/q_j$ in lowest terms.
After excluding the fixed coefficient primes, use
\[
 H(A,B)=AB\prod_j(2q_jA^m+p_jB^m).
\]
This form is squarefree and has degree $rm+2$: the nonzero distinct
shifts give disjoint simple root sets, disjoint also from $AB=0$.
By \eqref{eq:polyabc},
\[
 \log\rad\prod_j(2q_jA^m+p_jB^m)
              \ge(rm-2-\epsilon)h(M)-O(1),
\]
since $\log\rad(AB)\le2h(M)$.
The numerator product has logarithmic size at most $rm h(M)+O(1)$.
Its prime support of multiplicity at least two is bounded by the
difference between logarithmic size and logarithmic radical.
Rational cancellation affects only the excluded coefficient primes:
the displayed factors are coprime to $B$ outside that set.
\end{proof}

\section{The derivative height estimate}

Fix distinct integer shifts $c_1,\ldots,c_k$ and write
\[
 a=M^m,\quad x=a-z,\quad y=a+z,\quad t=xy=a^2-z^2,
\]
\[
 F(a,z)=G(x)G(y)=\prod_{i=1}^k((a+c_i)^2-z^2),\qquad
                         u=-\frac{2kaF}{F_a},
\]
where $F_a$ denotes the derivative with respect to $a$ while $z$ is
fixed. Retain $M,t,F,F_a\ne0$, and require
$2M^m+c_i+c_j\ne0$ for every nonzero sum $c_i+c_j$.
The fixed set $S$ contains the primes dividing $2k$ and
all primes at which distinct shifts cease to be distinct.

\begin{theorem}\label{thm:derivative-height}
Under ordinary rational $abc$, for every fixed $\epsilon>0$ with
$k-6-3\epsilon>0$,
\begin{equation}\label{eq:derivative-height}
 \log r^+_{\mathrm{odd},S}(u)
 \ge \frac{k-6-3\epsilon}{2}h(t)
                  -(4+\epsilon)h(M)-O(1).
\end{equation}
The constant is uniform in $M,z$.
\end{theorem}
\begin{proof}
Interchange $x,y$ if necessary so that $h(x)\ge h(y)$.
Start with the simple-prime support $T(x)$ of Lemma~\ref{lem:simple}.
Discard primes dividing the numerator or denominator of $M$, at a
cost at most $2h(M)$. At every remaining prime under consideration,
$2ka$ is a unit. Thus multiplying $-F/F_a$ by $2ka$ does not change
any valuation used below. Since $G(x)$ has positive valuation, $x$ is integral;
then $y=2a-x$ is integral too.

If $v_p(G(y))=0$, exactly one factor $L=x+c_i$ has valuation one.
The identity
\[
                         F_a=G'(x)G(y)+G(x)G'(y)
\]
shows that $F_a$ is a unit, and $v_p(u)=1$.
If $v_p(G(y))\ge2$, discard $p$ in the support $J(y)$.
This costs at most $(2+\epsilon)h(y)+O(1)$.

It remains to treat $v_p(G(x))=v_p(G(y))=1$.
There are unique factors $L=x+c_i$ and $R=y+c_j$ of valuation one.
Write $F=LRK$, where $K$ is a local unit. Differentiation gives
\[
                       F_a=(L+R)K+LRK_a.
\]
The derivative $K_a$ is integral. Since
$L+R=2a+c_i+c_j$, valuation one of this sum implies
$v_p(F_a)=1$ and again $v_p(u)=1$.
Otherwise $p^2\mid 2a+c_i+c_j$ locally.
The case $c_i+c_j=0$ is impossible here because $a$ is a unit.
Let $D(M)$ be the product over the distinct nonzero sums $c_i+c_j$.
Lemma~\ref{lem:power-multiple} bounds the remaining exceptional
support by $(2+\epsilon)h(M)+O(1)$.

Every surviving prime therefore occurs in
$r^+_{\mathrm{odd},S}(u)$. Consequently
\begin{align*}
 \log r^+_{\mathrm{odd},S}(u)
 &\ge(k-4-2\epsilon)h(x)-(2+\epsilon)h(y)\\
 &\hspace{12mm}-(4+\epsilon)h(M)-O(1)\\
 &\ge(k-6-3\epsilon)\max(h(x),h(y))
                  -(4+\epsilon)h(M)-O(1).
\end{align*}
Finally $h(t)\le h(x)+h(y)\le2\max(h(x),h(y))$.
\end{proof}

\section{One elliptic detector}

Retain the curve, point charts, subgroup and formal-comparison theorems
of \cite{Conditional,Unit}:
\[
 K=\Q(\sqrt2),\qquad E:y^2=x^3-(3-\sqrt2)^2x,\qquad
 2E(K)=\Z P,
\]
\[
 P=\left(\frac{57-28\sqrt2}{8},
          \frac{-660+417\sqrt2}{32}\right).
\]
The three point roles are $A=[n]P$, $H=[h]P$, $U=[j]P$.
Keep every original nonzero point, affine difference and quotient guard.
With $f=1/x$ and $\zeta=-x/y$, put
\[
 r=\frac{\zeta(U)}{\zeta(H)},\qquad
 r'=\overline{\frac{\zeta(U-A)}{\zeta(H-P)}},
\]
\[
 C_1=f(H),\quad C_2=f(U),\quad
 C_3=\overline{f(H-P)},\quad C_4=\overline{f(U-A)},\quad
 C_5=r-r',\quad C_6=r-X.
\]
For a fixed odd integer $B\ge3$, set
\begin{equation}\label{eq:scale}
 \beta=X^2+B,\quad T=\Tr x(A),\quad
 M=\beta^4+T^4,\quad a=M^3,\quad t=a^2-z^2=M^6-z^2,
\end{equation}
where $z\in\Q$ is the one retained slope parameter.
Choose the sixteen shifts
\[
                  \{c_i\}=\{-8,-7,\ldots,-1,1,\ldots,7,8\}.
\]
Define $F,u$ as in the preceding section. An expression without $a$
uses
\begin{equation}\label{eq:squared-expression}
 \mathcal A=M^6,\quad \mathcal Z=z^2,\quad
 \mathcal G(\mathcal A,\mathcal Z)=
 \prod_{j=1}^8\bigl((\mathcal A+j^2-\mathcal Z)^2-4j^2\mathcal A\bigr).
\end{equation}
Then $F=\mathcal G(a^2,z^2)$, $F_a=2a\mathcal G_{\mathcal A}$ and
\[
                         u=-16\frac{\mathcal G}{\mathcal G_{\mathcal A}}.
\]
The integer polynomial $\mathcal G_{\mathcal A}$ has content eight.
Thus the reduced displayed expression is
$u=-2\mathcal G/\mathcal H$, where
$\mathcal H=\mathcal G_{\mathcal A}/8$ is primitive.
Set
\begin{equation}\label{eq:balanced}
 D_i=\frac{C_i^2}{1+uC_i^2},\qquad
 W=\prod_{i=1}^6\left(1+\frac{D_i^2}{u}\right),\qquad
 V=\Tr(W^{-1}),\qquad
 w=\frac{u}{(1+u)^2},\qquad b=\frac{4w}{V^2}.
\end{equation}
All divisions, $M,t,F,F_a$, the two-slot product norm, and $V$ are
guarded nonzero. No nonzero condition is imposed on an individual $C_i$.

For the single pole block use two rational parameters $v,s$ and the
intrinsic equation of \cite{Intrinsic}:
\begin{equation}\label{eq:pole}
 \eta=7b,\qquad
 Y^2=P_\eta(v,s):=
 \eta(\eta+v^2)^2s^4-\eta(\eta-v^2)^2(1-\eta)^2.
\end{equation}
The root $Y$ is a finite quadratic generator, not an additional free
coordinate. Its cleared radicand norm is retained as a guard.

\begin{lemma}\label{lem:balanced}
Let $p$ be an odd prime unramified in $K$ and let $\nu=v_p(u)$ be odd.
At each place above $p$,
\[
                  v_p(D_i)=\max(2v_p(C_i),-\nu).
\]
If $\nu>0$, positivity of $v_p(D_i)$ is equivalent to positivity of
$v_p(C_i)$. If $\nu<0$, every factor of $W$ is congruent to one
modulo $p$, independently of poles among the $C_i$.
\end{lemma}
\begin{proof}
The two terms of $1+uC_i^2$ have distinct valuations, since
$\nu+2v_p(C_i)$ is odd when finite. The displayed identity follows.
For $\nu<0$, it gives
\[
                v_p(D_i^2/u)\ge-3\nu>0.
\]
Identically zero tests are included by assigning them valuation
$+\infty$.
\end{proof}

\begin{lemma}\label{lem:extract}
For every rational solution of the guarded source, if $X\ne n$ then
\[
              r^+_{\mathrm{odd},S}(u)\mid |\operatorname{num}(X-n)|.
\]
Here $S$ is the retained fixed exceptional set, containing all primes
at most $76176$.
\end{lemma}
\begin{proof}
Let $\nu=v_p(u)>0$ be odd. Then $v_p(w)=\nu$.
Every factor of $W$ has valuation
$\min(0,2v(D_i)-\nu)\le0$. Let $q\le0$ be the largest local
valuation of $W$ at a place above $p$. Then $v_p(V)\ge-q$.
The intrinsic pole equation excludes negative odd valuations of
$\eta=7b$; outside $S$ this gives
\[
                    1\le v_p(b)\le\nu+2q.
\]
At a place attaining $q$, each factor has valuation greater than
$-\nu/2$, which implies $v(D_i)>\nu/4>0$.
Lemma~\ref{lem:balanced} gives $v(C_i)>0$ for all six original tests.
The retained formal-comparison theorem now applies without a change:
one of $h,h-1$ is a local unit; the two truncated logarithmic ratios
are congruent to $j/h$ and $(j-n)/(h-1)$; their equality and the
target test imply $X\equiv n\pmod p$ and $p$-integrality of $X$.
This is the local-implication argument of \cite[Section 4]{Conditional}.
\end{proof}

\section{Height soundness}

Put $E_h=h(1:\beta:T)$ and take $\kappa=1001/1000$.
Ordinary rational $abc$ supplies this weak exponent and every use of
\eqref{eq:polyabc}. The retained binary-power lemma and the signed
mixed-power theorem \cite{Signed} give
\begin{equation}\label{eq:rootbounds}
 \alpha E_h\le h(M)+O(1),\qquad
 h(M)\le Qh(t)+O(1),\qquad
 \alpha=\frac{997}{1001},\quad Q=\frac{1002}{995}.
\end{equation}
The second estimate uses $t=M^6-z^2$; its denominator is
$6-5\kappa=199/200>0$.
The trace-height theorem also gives
\[
                     n^2\le A_0h(T)+J_0
\]
for fixed constants. These constants are independent of $B$.

Taking $\epsilon=1/100$ in Theorem~\ref{thm:derivative-height} gives
\[
 \gamma h(t)\le\log r^+_{\mathrm{odd},S}(u)+C h(M)+O(1),\qquad
             \gamma=\frac{997}{200},\quad C=\frac{401}{100}.
\]
The exact strict absorption inequality is
\begin{equation}\label{eq:strict}
 Q\left(C+\frac1{2\alpha}\right)
       =\frac{225373347}{49600750}<\frac{997}{200}.
\end{equation}
Its positive difference is $87545567/198403000$.
No further $h(M)$ loss occurs in Lemma~\ref{lem:extract}.

\begin{theorem}\label{thm:sound}
Under the stated arithmetic inputs and ordinary rational $abc$,
every guarded rational source solution has $X=n\in\Z$ for every
sufficiently large fixed odd $B$.
\end{theorem}
\begin{proof}
Suppose $X\ne n$. Lemma~\ref{lem:extract} and the input-height bound
$2h(X)\le h(\beta)\le E_h$ give
\[
 \log r^+_{\mathrm{odd},S}(u)
           \le\tfrac12 E_h+\log(1+|n|).
\]
Together with \eqref{eq:rootbounds}, this gives
\[
 \left(\gamma-Q\left(C+\frac1{2\alpha}\right)\right)h(t)
                          \le\log(1+|n|)+O(1).
\]
The coefficient is positive. The trace-height bound and
$h(T)\le E_h\le Qh(t)/\alpha+O(1)$ bound $n$, then $h(t)$ and
$E_h$, independently of $B$. This contradicts
$E_h\ge h(\beta)\ge\log B$ for sufficiently large $B$.
\end{proof}

\section{Prime preparation and completeness}

Fix $X=n\ne0$ and take $A=[n]P$. Then $M$ and $a$ are fixed positive
rationals, with $a>8$. Temporarily put $U=[n]H$. Use the retained
point and sign-branch avoidance argument to choose compact real and
fixed finite-place neighborhoods for $H$ and its required differences.
For each actual test that is not identically zero on this multiplication
family, avoid its zero set as well, so its valuation is fixed in the
chosen finite neighborhoods. Identically zero tests remain zero.
Such finite avoidance is permitted by the retained cyclic-point density
and congruence refinement. All original guards are retained.

The sixteen symmetric shifts make $F$ even in both $a$ and $z$.
In particular $F_a$ is divisible by $a$, and
\begin{equation}\label{eq:asymptotics}
 F\sim z^{32},\qquad F_a\sim-32a z^{30},\qquad
 u\sim z^2,\qquad w\sim\frac1{z^2}
\end{equation}
as $z$ tends to infinity at any fixed place.
Both leading expressions for $u,w$ are rational squares.
At the selected real character, choose $z>a+8$.
Every factor of $F$ is negative, so $F>0$, and
\[
 \frac{F_a}{F}=\sum_i\frac{2(a+c_i)}{(a+c_i)^2-z^2}<0.
\]
Thus $u>0$. The same inequalities may be required at every prepared
real character by increasing the real lower bound for $z$.

At a fixed finite place, make $z$ sufficiently large relative to the
already fixed test valuations. Then $D_i\sim1/u$ for every nonzero
test, and $D_i=0$ for an identically zero test. Hence
\[
                       W\longrightarrow1,\qquad V\longrightarrow2.
\]
The convergence is uniform on the chosen neighborhoods.
At two, \eqref{eq:asymptotics} makes $w$ a local square of arbitrarily
large positive even valuation, and $v_2(V)=1$. Consequently
\[
 b\in\Q_2^{\times2},\quad v_2(b)\ge4,\quad
 \eta=7b\in7\Q_2^{\times2}.
\]
At seven make $w$ a local square of positive valuation and $V$ a unit;
then $v_7(b)\ge0$. At the real places $W\to1$, $w\to0$, and
$b\to0$, so require $0<b<1/7$. Include all fixed primes dividing
the numerator or denominator of $a$, and the original exceptional
and avoidance primes, in this temporary preparation. It does not
enlarge the global exceptional set in the soundness argument.

Choose an integer $D>0$ with $D^2(a+c_i)\in\Z$ for all $i$.
Apply Sun's simultaneous prime theorem \cite[Theorem 4.1]{Sun} to
the thirty-two distinct coefficients $\pm D^2(a+c_i)$, with arbitrary
prescribed divisibility of its square-denominator parameter, sufficiently
large primes, and the real lower bound just chosen. It gives
\[
 p_i^\pm=N\pm D^2(a+c_i)L^2,\qquad
                         z=\frac{N}{D^2L^2},
\]
where all thirty-two $p_i^\pm$ are distinct positive primes.
Prescribed divisibility of $L$ makes $z$ arbitrarily large at the
fixed finite places: the common numerator is a unit at these places,
since otherwise the distinct selected primes would all be divisible
by the same fixed prime. The theorem supplies the same conjugate
reduction-order compatibility and finite avoidance as the retained
elliptic preparation; the exact common-progression and rational-contact
deductions are \cite[Section 2]{Dependencies}. The theorem permits
every fixed finite number of positions, so all thirty-two positions
are covered by its stated scope. Uniformity in their coefficients is
not required for completeness.

The exact factorization is
\begin{equation}\label{eq:primefactorization}
                    F=\frac{\prod_i p_i^-p_i^+}{D^{64}L^{64}}.
\end{equation}
At a selected prime, just one of the thirty-two linear factors has
valuation one. Its opposite factor is a unit, and therefore so is
the derivative of its quadratic factor, $2(a+c_i)$. Hence
$v_p(F_a)=0$ and $v_p(u)=v_p(w)=1$.
Refine the shared auxiliary index using the retained compatible
congruences, preserving the already selected finite and real neighborhoods.
At the successful place the first four original tests have valuation
at least two and the comparisons at least four. The transformed tests
therefore make $W\equiv1\pmod p$.
At the other place its reciprocal product is zero or one modulo $p$,
so $V$ is a unit. Thus $v_p(b)=1$.

We verify all other odd valuations, including new primes of $F_a$.
Away from the temporary preparation set, $a$ is a local unit.
If $z$ is integral and $p$ is not selected, \eqref{eq:primefactorization}
makes $F$ a unit, while $F_a$ is integral. Thus $v_p(u)\le0$.
If $z$ has negative valuation, the leading terms in
\eqref{eq:asymptotics} dominate, since the remaining coefficients are
integral, $a$ is a unit, and $p\ne2$. Hence
\[
 v_p(F)=32v_p(z),\quad v_p(F_a)=30v_p(z),\quad
                         v_p(u)=2v_p(z)
\]
is even. At the fixed numerator and denominator primes of $a$, the
prescribed dominance in \eqref{eq:asymptotics} gives
$v_p(u)=2v_p(z)$, again even. The same holds at the fixed primes
dividing $2k$, by making the prescribed dominance sufficiently strong.

It follows that every positive odd valuation of $u$ outside the fixed
preparation occurs at a selected prime. At every negative odd valuation,
Lemma~\ref{lem:balanced} makes both slots of $W$ congruent to one,
so $V\equiv2\pmod p$ and
$v_p(b)=v_p(w)=-v_p(u)>0$.
If $v_p(u)$ is even, then $v_p(w)$ is even too: when $v_p(u)=0$,
the denominator $(1+u)^2$ still contributes an even valuation.
Division by $V^2$ does not change parity. The fixed-place preparation
handles two and seven as above. Therefore $b$ has square reduced
denominator, and so does $\eta=7b$.

The actual input to the intrinsic pole satisfies
\[
 0<\eta<1,\qquad \eta\text{ has square reduced denominator},\qquad
 \eta\in7\Q_2^{\times2},\qquad v_2(\eta)\ge4.
\]
Burungale's theorem, in the intrinsic form \cite{Intrinsic}, supplies
the rational witnesses $v,s,Y$ in \eqref{eq:pole}.

\section{Guards and eleven coordinates}

At every prepared point-algebra character, $a>8$ and $z>a+8$.
Thus $F>0$, $-F_a>0$, $u>0$, $1+u>0$, and $w>0$.
Each new denominator $1+uC_i^2$ is positive in both real slots,
and each factor of $W$ is at least one. Therefore the new denominator,
product, trace and input norms are nonzero. The original point and
difference guard proof is unchanged.

The intrinsic witness pairs are Zariski dense in their parameter plane.
At any point-algebra character with value $\eta=c\ne0$, the
polynomial $P_c(v,s)$ is nonzero, with leading coefficient
$c(c+v^2)^2$ in $s$. Consequently the norm of its cleared radicand
is a nonzero polynomial in $v,s$. Density permits avoiding that norm
and all other retained pole guards simultaneously. No bound $c<1$
is required at the other characters, and no guard on $1-c$ is used.

The point algebra has rank $4^3=64$. Adjoining the single monic
quadratic pole root gives rank $128$. The guarded separator and norm
determinant of \cite{Unit} use one rational coordinate and force a
rational character of the entire algebra. Their separation proof
permits nonreal conjugate pole roots; nonzero radicand norms suffice.
The coordinate census is
\[
 1\text{ input}+6\text{ point coordinates}+1\text{ slope }z
                 +2\text{ pole parameters}+1\text{ encoder}=11.
\]
All displayed rational operations admit finite integer polynomial
specifications after clearing the guarded denominators.
In the retained coordinates, the height argument $t=M^6-z^2$ has
degree at most $48$. The displayed numerator and denominator of $u$
have degrees at most $768$ and $720$, respectively. These are local
expression bounds, not degrees of the final norm determinant.

\begin{theorem}\label{thm:eleven}
Under ordinary rational $abc$ and the retained proved arithmetic
inputs cited above, there is a computable family of integer polynomials
in eleven variables, including the input, whose input projection is
exactly $\Z$ for every sufficiently large fixed odd coefficient shift
$B$. The usual final factor $X$ includes zero. The sufficient shift
is not numerically specified, and no complete determinant-degree or
coefficient-height bound is asserted here.
\end{theorem}
\begin{proof}
Theorem~\ref{thm:sound} gives soundness. The prime preparation and
intrinsic pole construction give completeness for each nonzero integer
input, including every whole-character guard. The one-coordinate
encoder gives the stated polynomial and census.
\end{proof}

\begin{thebibliography}{9}\raggedright
\bibitem{RSZ}Z.~Rudnick, P.~Sarnak and A.~Zaharescu,
\emph{The distribution of spacings between the fractional parts of
$n^2\alpha$}, Inventiones Mathematicae 145 (2001), 37--57,
Appendix A.2, equation (A.2),
\url{https://www.math.tau.ac.il/~rudnick/papers/rsz.pdf}.
\bibitem{Conditional}\emph{Thirteen Variables from Direct Elliptic
Index Comparison}, repository manuscript,
\nolinkurl{papers/conditional_few_variables/direct_input13_20261110_v22.tex}.
\bibitem{Unit}\emph{A single norm determinant for the 25- and 13-variable
constructions}, repository manuscript,
\nolinkurl{papers/finite_algebra/unit_norm_20261114_v28.tex}.
\bibitem{Signed}\emph{A signed mixed-power height inequality},
companion manuscript,
\nolinkurl{papers/elliptic_heights/signed_mixed_power_height_20261117_v58.tex}.
\bibitem{Intrinsic}\emph{A two-coordinate pole test of coefficient
height four}, companion manuscript,
\nolinkurl{papers/pole_parameters/intrinsic_pole_20261117_v56.tex},
using A.~Burungale,
\emph{A Diophantine definition of $\Z$ in $\Q$},
arXiv:2611.16130v1, Theorem B,
\url{https://arxiv.org/abs/2611.16130}.
\bibitem{Sun}Z.-W.~Sun,
\emph{48 unknowns for Hilbert's tenth problem over the rationals},
arXiv:2610.20691v3, Theorem 4.1,
\url{https://arxiv.org/pdf/2610.20691v3}.
\bibitem{Dependencies}\emph{Arithmetic dependencies of the defining
polynomial}, repository manuscript,
\nolinkurl{papers/arithmetic_dependencies/arithmetic_dependencies_20261108_v8.tex}.
\end{thebibliography}
\end{document}
